\chapter{Urinalysis}

\section*{What Is This Test?}
Urinalysis is a routine examination of urine using visual inspection, reagent-strip chemistry, and microscopic examination of sediment when indicated. It can detect blood, protein, glucose, ketones, bilirubin, leukocyte esterase, nitrite, cells, casts, crystals, and microorganisms. It is a screening and diagnostic test, not a standalone measure of overall kidney function.

\section*{Alternative Names}
Complete Urinalysis; Routine Urine Test; Urine Routine and Microscopy; Urine R\&M; Urine Dipstick and Microscopy

\section*{Why It Is Ordered}
Evaluation of urinary symptoms, suspected urinary infection, hematuria, proteinuria, kidney disease, stones, diabetes, dehydration, liver or biliary disease, and selected systemic illnesses. It is also used to follow known kidney or urinary tract disorders and to assess abnormal findings from another test.

\section*{How This Test Is Performed}
A fresh, clean-catch midstream urine specimen is preferred. The laboratory examines color and appearance, performs reagent-strip testing, and examines centrifuged sediment microscopically when ordered or when screening findings warrant it. A urine culture is a separate test and should be collected using the laboratory's specified container and transport method.

\section*{Preparation, Risks, and Limitations}
No fasting is usually required. Avoid contamination with menstrual blood, stool, skin products, or excessive delay before delivery to the laboratory. Urine concentration, exercise, medications, vitamin C, collection technique, and storage can alter results. A dipstick abnormality often requires repeat testing or microscopy; urinalysis alone does not diagnose a urinary infection or determine the cause of kidney disease.

\section*{References}
\begin{enumerate}
  \item MedlinePlus. Urinalysis.
  \item Kidney Disease: Improving Global Outcomes. Clinical evaluation of kidney disease.
\end{enumerate}

\clearpage
\section*{Understanding Results and Follow-up}
The laboratory report should identify the specimen, method, units, reference interval or decision limit, and whether the result is positive, negative, abnormal, or inconclusive. Reference intervals vary with age, sex, pregnancy, specimen type, assay, and laboratory; a result outside a range is not automatically a diagnosis, and a result within a range does not always exclude disease. Discuss unexpected results with the ordering clinician, who can decide whether repeat or confirmatory testing, treatment, or referral is appropriate. Do not change medicines or delay urgent care based on an isolated result.


\section*{Regulatory Status and Authorization Date}
Regulatory status applies to the specific assay, instrument, manufacturer, and intended use rather than to the test name alone. No single approval date can be assigned to this test category. For a commercial assay, verify the current product in the relevant regulator's database: the U.S. Food and Drug Administration (FDA) clearance, approval, or authorization; India's Central Drugs Standard Control Organisation (CDSCO) licence or permission; the European Union In Vitro Diagnostic Regulation (EU IVDR) CE certificate issued through the applicable conformity-assessment route; or the United Kingdom Medicines and Healthcare products Regulatory Agency (MHRA) registration and UKCA/CE status. Laboratory-developed tests may not have an individual FDA, CDSCO, EU IVDR, or MHRA approval date.
\clearpage
